\documentclass[../dayan_kernel_manual.tex]{subfiles}

\begin{document}

\section{Host RPC 服务}\label{sec:kernel-060}

Dayan OS 主动请求的 opcode 1--8 必须由 Host 实现。Host 使用反方向的 RPC 发布新代码
并启动 kernel；响应仍走 Device--Host mailbox。

\subsection{服务列表}

\begin{dayantable}{RPC opcode}{tab:kernel-060}{X[13,l] X[30,l] X[23,l] X[34,l]}
  opcode & 名称 & 方向 & 用途 \\
  1 & \code{RPC_WRITE} & Device $\rightarrow$ Host & 输出一段 proc 数据。 \\
  2 & \code{RPC_CLEAR} & Device $\rightarrow$ Host & 要求输出回调刷新指定 proc 的下游输出流。 \\
  3 & \code{RPC_HEAP_ACQUIRE} & Device $\rightarrow$ Host & 为一个 hart 分配私有堆 region。 \\
  4 & \code{RPC_HEAP_RELEASE} & Device $\rightarrow$ Host & 归还此前分配的私有堆 region。 \\
  5 & \code{RPC_TRAP_REPORT} & Device $\rightarrow$ Host & 上报同步异常现场。 \\
  6 & \code{RPC_TENSOR_INTERRUPT_REPORT} & Device $\rightarrow$ Host & 上报张量机器外部中断快照。 \\
  7 & \code{RPC_PROC_EXIT} & Device $\rightarrow$ Host & 上报 Host root proc 的退出值。 \\
  8 & \code{RPC_PROFILE_FLUSH} & Device $\rightarrow$ Host & 上报一个 Hart 的满 profile FIFO。 \\
  \code{0x0100} & \code{RPC_PROC_CREATE} & Host $\rightarrow$ Device & 在目标 hart 创建 kernel proc。 \\
  \code{0x0101} & \code{RPC_FENCEI} & Host $\rightarrow$ Device & 在目标 hart 执行一次 \code{fence.i}。 \\
  \code{0x0102} & \code{RPC_PROC_PS} & Host $\rightarrow$ Device & 同步目标 hart 的 proc 快照。 \\
  \code{0x0103} & \code{RPC_PROC_KILL_ALL} & Host $\rightarrow$ Device & 终止目标 hart 的所有非系统 proc。 \\
  \code{0x0104} & \code{RPC_PROFILE_CONFIG} & Host $\rightarrow$ Device & 开启或停止目标 hart profile。 \\
  \code{0x0105} & \code{RPC_DCACHE_FLUSH} & Host $\rightarrow$ Device & 失效目标 hart 的一段 D-cache。 \\
  \code{0x8100} & \code{RPC_PROC_CREATE_RESPONSE} & Device $\rightarrow$ Host & 确认 create 请求已处理。 \\
  \code{0x8101} & \code{RPC_FENCEI_RESPONSE} & Device $\rightarrow$ Host & 确认本 hart 的 I-cache 刷新完成。 \\
  \code{0x8102} & \code{RPC_PROC_PS_RESPONSE} & Device $\rightarrow$ Host & proc 快照完成。 \\
  \code{0x8103} & \code{RPC_PROC_KILL_ALL_RESPONSE} & Device $\rightarrow$ Host & kill-all 请求完成。 \\
  \code{0x8104} & \code{RPC_PROFILE_CONFIG_RESPONSE} & Device $\rightarrow$ Host & profile 配置完成。 \\
  \code{0x8105} & \code{RPC_DCACHE_FLUSH_RESPONSE} & Device $\rightarrow$ Host & D-cache 范围失效完成。 \\
\end{dayantable}

Service 的 output、trap 和 proc-exit callback 分别接收 WRITE/CLEAR、TRAP REPORT 和
root 退出事件。堆服务使用 HAL 的普通 Device RAM allocation，并按 Cluster、Core 和
hart owner 记录所有权。

\subsection{通用 packet}

WRITE、CLEAR、HEAP ACQUIRE、HEAP RELEASE、PROC EXIT 和 PROFILE FLUSH 的 command argument 指向 40-byte
\code{rpc_packet}：

\begin{dayantable}{\code{rpc_packet}}{tab:kernel-061}{X[17,l] X[29,l] X[18,l] X[36,l]}
  offset & 字段 & 类型 & 含义 \\
  0 & \code{request_addr} & \code{uint64_t} & 请求数据的 Device 地址。 \\
  8 & \code{request_size} & \code{uint64_t} & 请求长度。 \\
  16 & \code{response_addr} & \code{uint64_t} & 响应区的 Device 地址。 \\
  24 & \code{response_size} & \code{uint64_t} & 响应区长度。 \\
  32 & \code{status} & \code{int64_t} & Host 写回的结果，负值表示失败。 \\
\end{dayantable}

packet、request 与 response 分别从独占的 64-byte 对齐 slot 开始，结构的 wire 长度仍按
表中定义。Host 校验三个地址的 64-byte 对齐、固定长度和 Device 地址范围。处理完成后
先写 response，最后写 packet 的 \code{status}；清除 outbound doorbell 后不得再修改
这些区域。Device 随后 invalidate Host 写过的 slot，再把 response 复制回调用者对象。

\subsection{输出服务}

WRITE 和 CLEAR 的 request 都是 24-byte \code{rpc_print_request}：

\begin{dayantable}{\code{rpc_print_request}}{tab:kernel-062}{X[18,l] X[28,l] X[18,l] X[36,l]}
  offset & 字段 & 类型 & 含义 \\
  0 & \code{hart_id} & \code{uint32_t} & 来源 hart，必须与 mailbox 一致。 \\
  4 & \code{proc_id} & \code{uint32_t} & 输出所属 proc。 \\
  8 & \code{data_addr} & \code{uint64_t} & WRITE 数据地址；CLEAR 为 0。 \\
  16 & \code{data_size} & \code{uint64_t} & WRITE 数据长度；CLEAR 为 0。 \\
\end{dayantable}

Host 按 Cluster/Core/Host-root-proc 区分输出流；子 Thread 的独立 Device printf buffer
也归并到所属 root ID，因而无需向 Host 发布子 Thread 生命周期。当前 Service 将单次 WRITE 限制为 1 MiB，并在读取前
检查整个数据范围位于 Device RAM。Server 根据正在执行的 proc 将数据放入对应 Runtime
Session 的队列；Runtime 在 Event 查询、等待或 Stream 同步后通过
\code{NPU_IPC_OUTPUT_READ} 取回数据，并在应用线程调用 output callback。因此算子输出
属于发起调用的应用，而不是 \code{dayan_server} 的标准输出。

\subsection{堆服务}

\begin{dayantable}{堆 RPC 数据}{tab:kernel-063}{X[29,l] X[18,l] X[53,l]}
  类型 & 大小 & 字段 \\
  \code{rpc_heap_request} & 8 & 请求的最小字节数。 \\
  \code{rpc_heap_response} & 16 & Host 返回的地址和实际大小。 \\
  \code{rpc_heap_release} & 16 & 待归还的地址和大小。 \\
\end{dayantable}

Host 把 ACQUIRE 大小向 64-byte 对齐，并返回不小于请求值的普通 RAM region。RELEASE
只有在地址、大小、类型和 hart owner 都匹配现有 allocation 时成功。

\subsection{异常报告}

TRAP REPORT 不使用 \code{rpc_packet}；command argument 直接指向 328-byte
\code{trap_record}。记录包含 hart/thread 标识、thread 状态、五个 machine/F CSR、
32 个整数寄存器，以及所属 Host root proc、parent proc 和 launch cookie。Host 校验
来源 hart 后调用 trap callback。

\subsection{创建 proc}

Host 把 224-byte \code{rpc_proc_create_request} 放在独占 cache line 的普通 Device RAM，
并保持到收到创建回复、读回 \code{RUNNING} 或 \code{FAILED}：

\begin{dayantable}{proc create 描述符}{tab:kernel-064}{X[26,l] X[18,l] X[56,l]}
  字段组 & offset & 含义 \\
  state & 0--7 & state 和 flags。 \\
  launch memory & 8--47 & entry、stack size、shared base/size、printf buffer size。 \\
  arguments & 48--119 & argc、reserved 和 8 个 64-bit 参数槽。 \\
  task resources & 120--127 & task mode 和最大子 Thread 数。 \\
  result & 128--151 & status、proc ID、hart ID、cookie。 \\
  logical task & 152--179 & 三维 task ID、三维 task dim 和一个 reserved 字段。 \\
  kernel resources & 180--207 & 32-bit Named Resource Context、URAM base/size 和 64-bit URAM/TCSR Context。 \\
  proc name & 208--223 & 固定 16-byte root proc 名称。 \\
\end{dayantable}

Host Service 校验启动参数，Dayan OS 只检查固件的静态容量，从本 Hart 的 slot 池取得
root 和子 slot，并把所有启动参数复制到私有对象。Device 将状态改为 \code{RUNNING}，
完整 flush 描述符后提交 root token。接收线程同步回复 opcode
\code{0x8100} 与原描述符地址、清除 MSIP 后才切换，新 kernel 不会阻塞创建回复。
该响应只表示创建完成，Host 读回后即可释放描述符。退出结果通过 \code{RPC_PROC_EXIT}
单独发送，Device 不保留创建描述符指针。

stack size 范围为 1--64 KiB，printf buffer 不超过 4 KiB，参数数量不超过 8。
root 与最多 15 个子 Thread 共用 16 个算子 slot；init、RPC 与 profile proc 不计入。
名称必须非空且在 16-byte 字段内以 NUL 结尾，所以有效名称最多为 15 bytes。
shared 区域为空时 base/size 同时为 0；非空时必须位于 coherent window 且 16-byte 对齐。
URAM 区域为空时 base/size 同时为 0；非空时必须落在请求来源 hart 的 URAM window，
base 和 size 都按 32 KiB 对齐。描述符还携带 64-bit Context1 和 22-bit Context2；后者
编码 Local/Shared Named Resource 的物理 base 和逻辑 count。最大子 Thread 数当前不超过
15。

\subsection{退出通知与管理 RPC}

每个 proc 退出都会先尝试用 WRITE 提交剩余输出；缓冲为空时不发送 WRITE，正常退出也不
额外发送 CLEAR。只有
\code{RPC_PROC_CREATE} 创建的 root 再发送 opcode 7；32-byte
\code{rpc_proc_exit_request} 含 hart/proc ID、必须为 0 的 parent ID、
root flag、\code{int64_t result} 和 launch cookie。Service 的 proc-exit callback 用
Cluster/Core/proc/cookie 定位 Host 资源。子 Thread 的创建和退出完全在 Device family
内完成，不发送该 RPC。自然退出的 root 在请求入队后停止执行，但其 slot 和承载
transaction 的栈会保留到 Host 清除 doorbell；随后由 send Thread 完成 descriptor cleanup
并回收 slot。

张量机器外部中断使用 opcode 6（\code{RPC_TENSOR_INTERRUPT_REPORT}）。command argument
直接指向 472-byte \code{tensor_interrupt_record}，包含 Hart/Thread、当前 Tensor
Context、中断控制值、累计计数、最近 16 条记录和软件丢弃计数。Service 校验
来源 Hart 与 recent count 后调用 tensor-interrupt callback；该报告不改变 Proc 生命周期。

\code{RPC_PROC_PS} 的参数指向 1864-byte 快照。40-byte header 含 hart、record 数、
全局空闲 slot、reserved slot 和 status，后接最多 19 条 96-byte record：三个系统
proc（receive、send、profile）加最多 16 个算子 proc。record 包含 proc/parent ID、thread state、root/child/zombie/
kill-wait flags、直属子计数、slot、Host root ID、logical task、返回值、PC 和 cookie。Host Service 通过
\code{npu_service_rpc_ps(device,target,...)} 显式指定 Cluster/Core。

\code{RPC_PROC_KILL_ALL} 使用 40-byte result，返回 killed root/child 数、已经释放的
slot 数和仍处于 Produce 延迟回收的数量。系统 proc 永不被终止；每个被终止 root 以
\code{-ECANCELED} 发送一次 root exit callback。这里的
\code{waiting_produce} 是尚未返回的硬件调度 token 数，也会计入 yield 提交的无依赖
Wait。Service 对应接口为 \code{npu_service_rpc_kill_all(device,target,...)}。

\subsection{proc profile 服务}

Host 用 56-byte \code{rpc_profile_config} 对每个目标 Hart START/STOP。描述符包含
command、hart、FIFO 地址/容量和 64-bit generation；STOP response 还返回未满尾批次的
字节数。Service 默认分配一个 16 MiB 的 Device-scope allocation，并按已启动 Hart
均分独立 FIFO slice，因此原始记录无需携带 Hart ID。每 Hart slice 至少容纳 19 组
24-byte 定义和 8-byte 初始状态，保证 START 可报告三个系统 proc 和最多 16 个算子 proc。

满缓冲请求为 32-byte \code{rpc_profile_flush_request}，包含 hart、地址、count、capacity
和 generation；24-byte response 包含 status 与 \code{continue_profile}。Host 必须先
读取 FIFO，再写 response，最后完成 packet。回调或队列无法继续接收时应返回
\code{continue_profile=0}，Device 随即关闭该 Hart 的记录。

\subsection{发布新代码}

Host 在 BL1 启动前为每个 Cluster 预留 16 MiB kernel arena，按 ELF 对齐要求单调
分配模块所有加载节段，包含 \code{.bss} 零初始化。卸载或上传失败不复用已写入地址；
arena 用尽返回 \code{-ENOMEM}，Service 关闭时统一释放。新模块因此无需 D-cache RPC。

Host 上传后向该 Cluster 所有已启动 Hart 并行发送 \code{RPC_FENCEI}，全部完成后
发布入口。argument 是不解引用的模块地址 cookie，旧 \code{RPC_ICACHE_FLUSH} 名称
保留为相同 opcode 的别名。loader 互斥锁覆盖分配、上传和发布；模块复用无需上传或 fence。
Profile START、STOP 及失败回滚使用同一个跨核并发 helper，同一 Hart 的 RPC 仍串行。

\code{RPC_PROC_CREATE} 本身不再执行 \code{fence.i}。create 描述符的请求字段由 Host
保持不变，因此 Device invalidate 后直接通过 \code{request->field} 校验和初始化；只有
result 字段与 state 由 Device 更新。

所有 Host-owned 标量控制 descriptor 从 64-byte 对齐地址开始，allocation 大小也按
64-byte 向上取整，opcode 唯一确定结构体布局和长度。Device 在解析前 invalidate，
请求字段保持只读；必须写入的结果和 state 在响应或 Proc family 终态完成，并在通知 Host
前 clean。Host 确认响应或终态后才能释放 descriptor。

\end{document}
